% !TeX root = ../navier-stokes-manuscript.tex
% ============================================================
% 2.1 Vortex Stretching
% ============================================================

The momentum equation has four pieces:
\[
  \underbrace{\partial_t u}_{\text{local acceleration}}
  +
  \underbrace{(u\cdot\nabla)u}_{\text{nonlinear transport}}
  =
  \underbrace{-\nabla p}_{\text{pressure}}
  +
  \underbrace{\nu\Delta u}_{\text{viscous diffusion}}.
\]
Nonlinear transport carries the velocity through the velocity field and can sharpen its spatial variation as it moves.
Viscosity diffuses that variation.
The hope is simple: viscosity smooths the fluid faster than nonlinear transport can sharpen it.
If that remains true at every scale, the velocity stays smooth.

\subsection{Vortex Stretching}\label{sub:ThePerfectlyStretchingVortex}

In three dimensions, strain can sharpen the local rotation carried by the fluid.
Set \(\omega:=\nabla\times u\) and write the symmetric strain as
\[
  S:=\frac12\bigl(\nabla u+(\nabla u)^{\mathsf T}\bigr),
\]
Taking the curl removes pressure, and its stretching contribution is \(S\omega\).
When a positive direction of the strain aligns with the vorticity, the surrounding flow can extend the rotating fluid along its axis while viscosity continues to diffuse it.
Take one material vortex segment aligned with that direction.
If \(e\) is its axial direction and \(L(t)\) its length, then
\[
  Se=\sigma(t)e,
  \qquad
  \sigma(t)>0,
  \qquad
  \frac{\dot L(t)}{L(t)}
  =
  e\cdot Se
  =
  \sigma(t).
\]
The segment lengthens with the fluid; it is not pulled through a fixed tube drawn around it.

As \(L(t)\) grows, incompressibility determines what happens across the segment at the same time.
Let \(\mathcal V\) be its fixed material volume, set \(A(t):=\mathcal V/L(t)\), and write \(r_{\mathrm{eff}}(t):=\sqrt{A(t)/\pi}\).
Then
\[
  \nabla\cdot u=0,
  \qquad
  \frac{d}{dt}(A L)=0,
  \qquad
  \frac{\dot A}{A}=-\sigma(t),
  \qquad
  \frac{\dot r_{\mathrm{eff}}}{r_{\mathrm{eff}}}
  =
  -\frac{\sigma(t)}{2}.
\]

The axial strain also acts on the rotation carried by the segment.
With \(\omega:=\nabla\times u\),
\[
  \frac{D\omega}{Dt}
  =
  S\omega+\nu\Delta\omega,
  \qquad
  \frac{D}{Dt}
  :=
  \partial_t+u\cdot\nabla,
\]
and when \(\omega=|\omega|e\),
\[
  S\omega
  =
  \sigma(t)\omega,
  \qquad
  \frac12\frac{D|\omega|^2}{Dt}
  =
  \sigma(t)|\omega|^2
  +
  \nu\,\omega\cdot\Delta\omega.
\]

On an interval where the right-hand side stays positive, the rotation strengthens.
If \(|u|\leq U_0\) on the material segment \(\Omega_{\mathrm{seg}}(t)\), then
\[
  E_{\mathrm{seg}}(t)
  :=
  \frac12\int_{\Omega_{\mathrm{seg}}(t)}|u(x,t)|^2\,dx
  \leq
  \frac12U_0^2\mathcal V
  <
  \infty,
  \qquad
  \left|\frac12\bigl(\nabla u-(\nabla u)^{\mathsf T}\bigr)\right|_{\mathrm F}
  =
  \frac{|\omega|}{\sqrt2}
  \leq
  |\nabla u|_{\mathrm F}.
\]
The segment can carry finite kinetic energy while the rotating part of its velocity gradient grows.
The concentration is showing itself through the falling width, not through a forced increase in energy.

\divider{\NavierStokes{} Scaling}

The narrowing has made scale the variable that tests the original viscosity hope.
Fix \(t_0\) and set \(\ell:=r_{\mathrm{eff}}(t_0)\).
For \(\lambda>1\), define
\[
  u_\lambda(x,t)
  :=
  \lambda u(\lambda x,\lambda^2t),
  \qquad
  p_\lambda(x,t)
  :=
  \lambda^2p(\lambda x,\lambda^2t).
\]
At time \(\lambda^{-2}t_0\), the vortex in \(u_\lambda\) has width \(\lambda^{-1}\ell\).
This produces another solution at a different scale, not a later state of the original vortex.
Its four momentum terms are
\[
\begin{aligned}
  \partial_tu_\lambda(x,t)
  &=
  \lambda^3(\partial_tu)(\lambda x,\lambda^2t),
  \\
  (u_\lambda\cdot\nabla)u_\lambda(x,t)
  &=
  \lambda^3\bigl((u\cdot\nabla)u\bigr)(\lambda x,\lambda^2t),
  \\
  \nabla p_\lambda(x,t)
  &=
  \lambda^3(\nabla p)(\lambda x,\lambda^2t),
  \\
  \nu\Delta u_\lambda(x,t)
  &=
  \lambda^3\nu(\Delta u)(\lambda x,\lambda^2t).
\end{aligned}
\]
The finer scale strengthens viscosity and nonlinear transport by the same \(\lambda^3\) factor; pressure and local acceleration follow with them.
Going smaller has not given viscosity a scaling advantage.

\divider{Critical Height}

The equation has kept its balance while the norms used to measure the velocity change with scale.
Let \(\Lambda:=(-\Delta)^{1/2}\).
A change of variables gives
\[
  \norm{u_\lambda(t)}_{\dot H^s}^2
  =
  \lambda^{2s-1}
  \norm{u(\lambda^2t)}_{\dot H^s}^2.
\]
The exponent vanishes at \(s=\tfrac12\).
Set
\[
  H(t)
  :=
  \frac12\norm{u(t)}_{\dot H^{1/2}}^2
  =
  \frac12\norm{\Lambda^{1/2}u(t)}_2^2.
\]
With \(H_\lambda\) denoting the critical height of \(u_\lambda\),
\[
  \norm{u_\lambda(t)}_2^2
  =
  \lambda^{-1}\norm{u(\lambda^2t)}_2^2,
  \qquad
  H_\lambda(t)=H(\lambda^2t),
  \qquad
  \norm{\nabla u_\lambda(t)}_2^2
  =
  \lambda\norm{\nabla u(\lambda^2t)}_2^2.
\]
Kinetic energy falls as first derivatives rise.
The critical height remains fixed.
At \(\dot H^{1/2}\), the concentration that kinetic energy can miss is measured without the scale growth of the first-derivative norm.

\divider{Nonlinear Production versus Viscous Dissipation}

Along the original solution, \(H(t)\) can still rise or fall.
Apply \(\Lambda^{1/2}\) to the momentum equation and pair with \(\Lambda^{1/2}u\).
The signed nonlinear production is
\[
  \mathcal P_{1/2}(t)
  :=
  -\operatorname{Re}
  \pair
    {\Lambda^{1/2}u(t)}
    {\Lambda^{1/2}\bigl((u(t)\cdot\nabla)u(t)\bigr)}_{L^2}.
\]
The pressure pairing vanishes because \(\nabla\cdot u=0\), while the Laplacian contributes \(-\nu\norm{\Lambda^{3/2}u}_2^2\).
The critical height obeys
\[
  H'(t)
  +
  \nu\norm{\Lambda^{3/2}u(t)}_2^2
  =
  \mathcal P_{1/2}(t).
\]
It rises exactly when nonlinear production exceeds viscous dissipation.
The two contributions scale as
\[
  \mathcal P_{1/2}[u_\lambda](t)
  =
  \lambda^2\mathcal P_{1/2}[u](\lambda^2t),
  \qquad
  \nu\norm{\Lambda^{3/2}u_\lambda(t)}_2^2
  =
  \lambda^2\nu\norm{\Lambda^{3/2}u(\lambda^2t)}_2^2.
\]
Production and dissipation remain matched in scaling.
Their common \(\lambda^2\) factor is the inverse of the time contraction \(t\mapsto\lambda^{-2}t\).

\divider{The Heat Clock}

The viscous side of that balance carries a width-dependent comparison clock.
For viscosity alone, the relation between width and time is visible in the linear heat equation \(v_t=\nu\Delta v\):
\[
  \widehat v(\xi,t+\delta t)
  =
  e^{-\nu|\xi|^2\delta t}\widehat v(\xi,t).
\]
A scale-localized structure of width \(r\) occupies frequencies \(|\xi|\simeq r^{-1}\).
Its viscous change becomes order one when \(\nu|\xi|^2\delta t\simeq1\), giving the comparison time
\[
  \tau_\nu(r)
  :=
  \frac{r^2}{\nu}.
\]
At the finer width,
\[
  \tau_\nu(\lambda^{-1}r)
  =
  \lambda^{-2}\tau_\nu(r),
\]
which is the same time contraction carried by the full \NavierStokes{} rescaling.
Each decrease in width shortens the heat clock by the square of that decrease.

\divider{The Zeno Cascade}

Repeating the decrease in width repeats the decrease in time.
For the geometric sequence
\[
  r_n:=2^{-n}r_0,
  \qquad
  \tau_n:=\frac{r_n^2}{\nu},
\]
the heat clocks satisfy
\[
  \tau_n
  =
  \frac{r_0^2}{\nu}4^{-n},
  \qquad
  \sum_{n=0}^{\infty}\tau_n
  =
  \frac{4r_0^2}{3\nu}
  <
  \infty.
\]
Their summability shows that the clocks alone do not prevent infinitely many scale transitions from accumulating at finite time.
For them to describe a candidate singularity, one \NavierStokes{} solution would have to pass through widths
\[
  r_0>r_1>r_2>\cdots\longrightarrow0
\]
at times
\[
  t_0<t_1<t_2<\cdots\uparrow T_*<\infty,
  \qquad
  t_{n+1}-t_n\asymp\tau_n,
\]
with comparison constants independent of \(n\).
The remaining time then satisfies
\[
  T_*-t_n
  \asymp
  \sum_{k=n}^{\infty}\tau_k
  =
  \frac{4r_n^2}{3\nu}.
\]
Under that candidate schedule, each new width is formed inside an interval tending to zero at order \(\tau_n\).
The spatial sequence has become a demand on the rate at which the same velocity field can change.
At a fixed spatial coordinate, that demand begins with the velocity, continues through its local acceleration and jerk, and does not stop there.
It generates the temporal Tower.
